% GENERATED FILE. Edit canonical lessons, then run npm run generate:books.
% canonical-source-hash: fnv1a64:052077a6a09c0c73
% canonical-lessons: MR-C72-rakta, MR-C72-gudgha, MR-C72-doke, MR-C72-tap, MR-C72-davakhana

\chapter{Blood, Knee, Head, Fever, Clinic}
\label{ch:mr-blood-knee-head-fever-clinic}
\hlchaptermodality{72}

\begin{chapteropening}
\textbf{By the end of this chapter:} \emph{I can say blood, a knee, the head, a fever and a clinic.}

Complete the last lesson of chapter 72: I can say blood, a knee, the head, a fever and a clinic.
\end{chapteropening}

\section[\texorpdfstring{\mr{रक्त} (rakta)}{रक्त (rakta)}]{\mr{रक्त} (rakta) — blood}

\label{lesson:MR-C72-rakta}



\begin{quote}
Before the new one: say the Marathi for the neck, then the Marathi for the face.
\end{quote}

\subsection*{You'll want to know: \mr{रक्त}}
\textbf{\mr{रक्त}} — \emph{rakta} — \textquotedblleft{}blood\textquotedblright{}.

\begin{cousinweb}
\textbf{\mr{रक्त}} (\emph{rakta}) is Sanskrit \emph{rakta}, \textquotedblleft{}coloured, red; blood\textquotedblright{}, taken over whole as a learned loan, a tatsama (CDIAL 10539). It is the past participle of \emph{rañj}, \textquotedblleft{}to colour\textquotedblright{}.
\end{cousinweb}

\begin{grammarlens}[title={What you've built}]
One more word for this chapter.
\end{grammarlens}

\subsection*{Your turn}
\begin{itemize}
\raggedright
  \item \emph{Say it:} \emph{rakta}
  \item \emph{Say it:} \emph{rakta}, once more
  \item \emph{Say it:} say \emph{rakta}
  \item \emph{Recall:} say the Marathi for the neck, then the Marathi for the face, then say \emph{rakta} again
\end{itemize}

\begin{tcolorbox}[breakable,colback=teal!4,colframe=teal!35!black,title={Before you move on}]
What does \mr{रक्त} mean? (\textquotedblleft{}Blood\textquotedblright{}.) Say it once more.
\end{tcolorbox}
\section[\texorpdfstring{\mr{गुडघा} (guḍghā)}{गुडघा (guḍghā)}]{\mr{गुडघा} (guḍghā) — a knee}

\label{lesson:MR-C72-gudgha}



\begin{quote}
Before the new one: say the Marathi for the face, then the Marathi for blood.
\end{quote}

\subsection*{You'll want to know: \mr{गुडघा}}
\textbf{\mr{गुडघा}} — \emph{guḍghā} — \textquotedblleft{}a knee\textquotedblright{}.

\begin{cousinweb}
The origin of \textbf{\mr{गुडघा}} (\emph{guḍghā}) is not settled. Neither CDIAL nor DEDR, searched by form and by meaning, has an entry for it.
\end{cousinweb}

\begin{grammarlens}[title={What you've built}]
One more word for this chapter.
\end{grammarlens}

\subsection*{Your turn}
\begin{itemize}
\raggedright
  \item \emph{Say it:} \emph{guḍghā}
  \item \emph{Say it:} \emph{guḍghā}, once more
  \item \emph{Say it:} say \emph{guḍghā}
  \item \emph{Recall:} say the Marathi for the face, then the Marathi for blood, then say \emph{guḍghā} again
\end{itemize}

\begin{tcolorbox}[breakable,colback=teal!4,colframe=teal!35!black,title={Before you move on}]
What does \mr{गुडघा} mean? (\textquotedblleft{}A knee\textquotedblright{}.) Say it once more.
\end{tcolorbox}
\section[\texorpdfstring{\mr{डोके} (ḍoke)}{डोके (ḍoke)}]{\mr{डोके} (ḍoke) — the head}

\label{lesson:MR-C72-doke}



\begin{quote}
Before the new one: say the Marathi for blood, then the Marathi for a knee.
\end{quote}

\subsection*{You'll want to know: \mr{डोके}}
\textbf{\mr{डोके}} — \emph{ḍoke} — \textquotedblleft{}the head\textquotedblright{}.

\begin{cousinweb}
\textbf{\mr{डोके}} (\emph{ḍoke}) goes back to \emph{*ḍōkka}, \textquotedblleft{}head\textquotedblright{}, a form CDIAL reconstructs beside \emph{*ḍōka} (CDIAL 5566).
\end{cousinweb}

\begin{grammarlens}[title={What you've built}]
One more word for this chapter.
\end{grammarlens}

\subsection*{Your turn}
\begin{itemize}
\raggedright
  \item \emph{Say it:} \emph{ḍoke}
  \item \emph{Say it:} \emph{ḍoke}, once more
  \item \emph{Say it:} say \emph{ḍoke}
  \item \emph{Recall:} say the Marathi for blood, then the Marathi for a knee, then say \emph{ḍoke} again
\end{itemize}

\begin{tcolorbox}[breakable,colback=teal!4,colframe=teal!35!black,title={Before you move on}]
What does \mr{डोके} mean? (\textquotedblleft{}The head\textquotedblright{}.) Say it once more.
\end{tcolorbox}
\section[\texorpdfstring{\mr{ताप} (tāp)}{ताप (tāp)}]{\mr{ताप} (tāp) — a fever}

\label{lesson:MR-C72-tap}



\begin{quote}
Before the new one: say the Marathi for a knee, then the Marathi for the head.
\end{quote}

\subsection*{You'll want to know: \mr{ताप}}
\textbf{\mr{ताप}} — \emph{tāp} — \textquotedblleft{}a fever\textquotedblright{}.

\begin{cousinweb}
\textbf{\mr{ताप}} (\emph{tāp}) is Sanskrit \emph{tāpa}, \textquotedblleft{}heat, glow\textquotedblright{}, also \textquotedblleft{}pain\textquotedblright{}, a tatsama (CDIAL 5767). CDIAL lists the inherited Marathi form as \emph{tāv}, \textquotedblleft{}heating to red heat\textquotedblright{}.
\end{cousinweb}

\begin{grammarlens}[title={What you've built}]
One more word for this chapter.
\end{grammarlens}

\subsection*{Your turn}
\begin{itemize}
\raggedright
  \item \emph{Say it:} \emph{tāp}
  \item \emph{Say it:} \emph{tāp}, once more
  \item \emph{Say it:} say \emph{tāp}
  \item \emph{Recall:} say the Marathi for a knee, then the Marathi for the head, then say \emph{tāp} again
\end{itemize}

\begin{tcolorbox}[breakable,colback=teal!4,colframe=teal!35!black,title={Before you move on}]
What does \mr{ताप} mean? (\textquotedblleft{}A fever\textquotedblright{}.) Say it once more.
\end{tcolorbox}
\section[\texorpdfstring{\mr{दवाखाना} (davākhānā)}{दवाखाना (davākhānā)}]{\mr{दवाखाना} (davākhānā) — a clinic}

\label{lesson:MR-C72-davakhana}



\begin{quote}
Before the new one: say the Marathi for the head, then the Marathi for a fever.
\end{quote}

\subsection*{You'll want to know: \mr{दवाखाना}}
\textbf{\mr{दवाखाना}} — \emph{davākhānā} — \textquotedblleft{}a clinic\textquotedblright{}.

\begin{cousinweb}
\textbf{\mr{दवाखाना}} (\emph{davākhānā}) is a Persian compound: \emph{davā}, \textquotedblleft{}medicine\textquotedblright{}, from Arabic \emph{dawāʾ}, \textquotedblleft{}remedy\textquotedblright{}, whose root is \emph{d-w-y} (Lane), and Persian \emph{khāna}, \textquotedblleft{}house\textquotedblright{}.
\end{cousinweb}

\begin{grammarlens}[title={What you've built}]
One more word for this chapter.
\end{grammarlens}

\subsection*{Your turn}
\begin{itemize}
\raggedright
  \item \emph{Say it:} \emph{davākhānā}
  \item \emph{Say it:} \emph{davākhānā}, once more
  \item \emph{Say it:} say \emph{davākhānā}
  \item \emph{Recall:} say the Marathi for the head, then the Marathi for a fever, then say \emph{davākhānā} again
\end{itemize}

\begin{tcolorbox}[breakable,colback=teal!4,colframe=teal!35!black,title={Before you move on}]
What does \mr{दवाखाना} mean? (\textquotedblleft{}A clinic\textquotedblright{}.) Say it once more.
\end{tcolorbox}
